\songtitle{Tres Ciutats, Una Mateixa Freqüència}

\tag*{2026}\tag{ebm}\tag{catalan-lyrics}\tag{ai-portrait}
\begin{notes}
Lyrics by SUNO based on a summary by Claude by examining characters and images by Stable Diffusion prompted by Carlos Thompson.
\end{notes}
\begin{style}
dark industrial club, EBM-tech basslines, rapid-fire pop chorus, Catalan storytelling, bilingual verses, whispered female mezzo vocals, metallic percussion, distorted synth stabs, four-on-the-floor kick, sidechain pumping, tape saturation, plate reverb, stereo delay throws, electric guitar riffs, acoustic guitar strums, analog synth pads, layered keyboards, dynamic breakdowns, build-and-drop arrangement, propulsive 146 BPM
\end{style}
\begin{lyrics}
\songpart{Verse 1 — Ashlee Hill}
Des de Iowa fins Chicago, un vol sense retorn,\\
cabells vermells, ulleres, auriculars al coll.\\
Entre pòsters de bandes i cables escampats,\\
el món sencer li cap dins dels seus auriculars.\\
No busca companyia, no demana permís,\\
té la seva pròpia lliga, el seu propi país.\\
Però què passa quan el codi, tan ben programat,\\
no té resposta per allò que no ha esperat?

\songpart{Pre-Chorus}
No és només so, no és només joc,
és una armadura, és un foc.\\
Però fora hi ha un món que no entén\\
de pistes, ritmes ni de béns.

\songpart{Chorus}
Tres ciutats, una mateixa freqüència,\\
tres camins que busquen la seva essència.\\
Ashlee, Nana, Silvia — noms sense edat,\\
cada una porta el seu propi combat.\\
No és fugir, és trobar-se en el camí,\\
tres dones, tres veus, un mateix destí.

\songpart{Vers 2 — Adriana "Nana" Quispe}
Lima, Stockholm, Barcelona — mapes a la pell,\\
els cabells de dos colors com un senyal d'atzar.\\
A la biblioteca, amb jaqueta de cuir,\\
somriu com qui sap que no ha de demostrar.\\
El cognom dels seus avis, l'herència andina,\\
barreja amb l'escena punk, la vida clandestina.\\
Comunicadora, mediàtica, urbana,\\
no tria entre el saber i la gana de viure.

\songpart{Pre-Chorus}
No és només estil, no és només veu,
és la història que dansa al seu peu.\\
Tres continents dins d'un mateix cos,\\
tres llengües, tres mons, un sol compàs.

\songpart{Chorus}
Tres ciutats, una mateixa freqüència,\\
tres camins que busquen la seva essència.\\
Ashlee, Nana, Silvia — noms sense edat,\\
cada una porta el seu propi combat.\\
No és fugir, és trobar-se en el camí,\\
tres dones, tres veus, un mateix destí.

\songpart{Bridge}
L'Amèrica profunda, els Andes, la mar,\\
l'illa sarda que es nega a callar.\\
Totes porten dins un viatge ancestral,\\
una memòria que es nega a morir.\\
I en aquesta Barcelona que les veu passar,\\
les tres es troben sense saber-ho encara.

\songpart{Verse 3 — Silvia Lai}
Sardenya a la sang, illa dins del mar,\\
cabell fosc, piga al rostre, somriure solar.\\
Samarreta blanca, arracades d'or,\\
abraçada a un ós de peluix — sense por.\\
No es pren la vida massa seriosament,\\
però sap que l'alegria és també un instrument.\\
De l'illa als carrers de Barcelona,\\
porta la calidesa que les altres no sonen.

\songpart{Pre-Chorus}
No és només riure, no és només dolç,
és l'equilibri que sosté el seu pols.\\
Quan les altres dubten, ella fa un pas,\\
quan el món pesa, ella sap que no és un fracàs.

\songpart{Final Chorus}
Tres ciutats, una mateixa freqüència,\\
tres camins que busquen la seva essència.\\
Ashlee, Nana, Silvia — noms sense edat,\\
cada una porta el seu propi combat.\\
No és fugir, és trobar-se en el camí,\\
tres dones, tres veus, un mateix destí.

Del Mig-Oest als fiords, dels Andes al mar,\\
de l'illa sarda al carrer de Gràcia,\\
totes tres cerquen el que ja tenen dins:\\
la força de ser qui són, sense fingir.\\
I si el món no entén la seva lletra,\\
elles canten juntes — la seva pròpia orquestra.\\
Tres ciutats, una veu, un mateix respir,\\
tres històries que aprenen a compartir.
\end{lyrics}

